% source: https://github.com/pfradin/luadraw/discussions/221#discussioncomment-16034961 (pfradin, block 1)
\documentclass{standalone}
\usepackage[svgnames]{xcolor}
\usepackage{luadraw}
% https://tex.stackexchange.com/a/754350/322482
\begin{document}
    \begin{luadraw}{name=two_ellipse}
        -- local g = graph:new{window={-5,5,-3,6}} -- good!
        local g = graph:new{window={-10,10,-9,9}, size={16,14}}
        -------------------------------------
        
        local i,sin,cos,pi = cpx.I, math.sin,math.cos,math.pi
        -- ellipse 1 (E1) : center c1, half axes a1 (directed by the vector u1) and b1
        -- ellipse 2 (E2) : center c2, half axes a2 (directed by the vector u2) and b2
        local c1,a1,b1,u1 = -1,4,2,1+i/2 -- (E1)
        local c2,a2,b2,u2 -- (E2)
        local example = function(k) -- to choose (E2)
            if k == 1 then -- (E2) inside (E1)        
                c2,a2,b2,u2 = 0.5+i/2, 1.5, 1, 1-i 
            elseif k == 2 then -- (E2) outside (E1)
                c2,a2,b2,u2 = 2+4*i, 1.5, 1, 1-i
            elseif k == 3 then -- (E2) cuts (E1) in 2 points
                c2,a2,b2,u2 = -1+2*i, 2, 1, 1-i
            elseif k == 4 then -- (E2) cuts (E1) in 4 points
                c2,a2,b2,u2 = 2*i, 4, 1, 1-i 
            end
        end

        local ellipsetangence = function(k) -- to choose example
            example(k) -- we choose a example
            u1 = u1/cpx.abs(u1); u2 = u2/cpx.abs(u2)
            -- we will place ourselves in the coordinate system {c1,u1,i*u1}
            -- where the ellipse (E1) has the equation x^2/a1^2+y^2/b1^2=1
            local matrix = {c1,u1,i*u1}
            local invm = invmatrix(matrix)
            c2 = applymatrix(c2,invm) -- center of the ellipse (E2) in the new coordinate system
            u2 = applyLmatrix(u2,invm) -- major axis of the ellipse (E2) in the new coordinate system
            local alpha = cpx.arg(u2)*rad -- inclination of (E2) in the new coordinate system

            local p = function(t) -- parametrization of (E2)
                return c2+a2*cos(t)*u2+b2*sin(t)*i*u2
            end
            local dp = function(t) -- derivation
                return -a2*sin(t)*u2+b2*cos(t)*i*u2
            end
            local ps = function(z1,z2) -- dot product related to (E1)
                return z1.re*z2.re/a1^2+z1.im*z2.im/b1^2
            end
            local N2 = function(z) -- squared norm
                return ps(z,z)
            end
            -- Let A=p(t) and u=dp(t), we seek t such that the equation N2(A+lambda*u)=1 
            -- has a single solution lambda, the discriminant must be zero.
            local T = solve(function(t) return ps(p(t),dp(t))^2-N2(dp(t))*(N2(p(t))-1) end, -pi, pi)
            if T == nil then T = {} end
            local sol = {}
            for _,t in ipairs(T) do
                local lambda = -ps(p(t),dp(t))/N2(dp(t)) -- solution double
                local A, u = p(t), dp(t)
                local B = A+lambda*u
                -- there are 4 common tangents; 
                -- you must select those for which the two centers are on the same side.
                if cpx.dot(A,i*u)*cpx.dot(A-c2,i*u) >=0 then
                    insert(sol,{A,B})
                end
            end
            -- now we draw
             g:Savematrix()  -- save the current matrix
            g:Composematrix(matrix)  -- compose matrix with the current graph matrix
            if #sol > 0 then
                local A,D,B,C = table.unpack(sol) -- first solution, the four points must be arranged in the correct order.
                local sens = -1;
                if cpx.det(B-A,D-A)<0 then sens = 1 end
                g:Dpath( {A,c2,B,a2,b2,sens,alpha,"ea",C,"l",0,D,a1,b1,sens,"ea","cl"}, "draw=none,fill=Pink");
                g:Lineoptions(nil,"blue",8); g:Dellipse(0,a1,b1); g:Dellipse(c2,a2,b2,alpha)
                g:Linecolor("ForestGreen")
                g:Dline(A,D); g:Dline(B,C)
                for k = 5, #sol-1, 2 do -- other tangents solutions?
                    g:Dline(sol[k],sol[k+1],"gray,dashed")
                end
                g:Dlabel(
                    "$(E_1)$",0,{},"$(E_2)$",c2,{},
                    "$A$", A, {pos="S",dist=0}, 
                    "$B$",B,{pos="N"}, "$C$",C,{}, 
                    "$D$",D,{pos="S"}
                )
            else
                g:Lineoptions(nil,"blue",8)
                g:Dellipse(0,a1,b1); g:Dellipse(c2,a2,b2,alpha)
                g:Dlabel(
                    "$(E_1)$",-a1,{pos="W"}, 
                    "$(E_2)$",c2+b2*i*u2,{pos="E"}, 
                    "no solution",-3*i,{pos="N"}
                )
            end
            g:Restorematrix() -- restore the initial matrix
        end

        -------------------------------
        g:Saveattr()
        g:Viewport(-10,0,-9,0) -- No.1
        g:Coordsystem(-5,5,-3,6)  -- this instruction modify the current matrix
        ellipsetangence(1)
        g:Restoreattr()

        g:Saveattr()
        g:Viewport(-10,0,0,9) -- No.2
        g:Coordsystem(-5,5,-3,6)
        ellipsetangence(2)
        g:Restoreattr()

        g:Saveattr()
        g:Viewport(0,10,-9,0) -- No.3
        g:Coordsystem(-5,5,-3,6)
        ellipsetangence(3)
        g:Restoreattr()

        g:Saveattr()
        g:Viewport(0,10,0,9) -- No.4
        g:Coordsystem(-5,5,-3,6)
        ellipsetangence(4)
        g:Restoreattr()

        g:Show()
    \end{luadraw}
\end{document}
